% ---
% Informações de dados para CAPA e FOLHA DE ROSTO
% ---
\titulo{[TO DO] Seu Título}
\autor{[TO DO] Nome Completo}
\local{Diamantina/MG}
\data{2023}
\orientador{Prof. Dr. Marcelo Ferreira Rego}
%\coorientador{Equipe \abnTeX}
\instituicao{%
  UNIVERSIDADE FEDERAL DOS VALES DO JEQUITINHONHA E MUCURI
  \par
  Bacharelado em Sistemas de Informação 
  }
\tipotrabalho{Monografia (Graduação)}

% O preambulo deve conter o tipo do trabalho, o objetivo, 
% o nome da instituição e a área de concentração 
\preambulo{Monografia apresentada ao Curso de Sistemas de Informação do Departamento de Computação da Universidade Federal dos Vales do Jequitinhonha e Mucuri como requisito parcial para obtenção do título de Bacharel em Sistemas de Informação.}
% ---


\pretextual

% ---
% Capa (Obrigatório)
% ---
\imprimircapa
% ---

% ---
% Folha de rosto (obrigatório)
% (o * indica que haverá a ficha bibliográfica)
% ---
\imprimirfolhaderosto
% ---


% ---
% Inserir a ficha bibliografica
% ---

% Isto é um exemplo de Ficha Catalográfica, ou ``Dados internacionais de
% catalogação-na-publicação''. Você pode utilizar este modelo como referência. 
% Porém, provavelmente a biblioteca da sua universidade lhe fornecerá um PDF
% com a ficha catalográfica definitiva após a defesa do trabalho. Quando estiver
% com o documento, salve-o como PDF no diretório do seu projeto e substitua todo
% o conteúdo de implementação deste arquivo pelo comando abaixo:
%
% \begin{fichacatalografica}
%     \includepdf{fig_ficha_catalografica.pdf}
% \end{fichacatalografica}
\begin{comment}
\begin{fichacatalografica}
	\vspace*{\fill}					% Posição vertical
	\hrule							% Linha horizontal
	\begin{center}					% Minipage Centralizado
	\begin{minipage}[c]{12.5cm}		% Largura
	
	\imprimirautor
	
	\hspace{0.5cm} \imprimirtitulo  / \imprimirautor. --
	\imprimirlocal, \imprimirdata-
	
	\hspace{0.5cm} \pageref{LastPage} p. : il. (algumas color.) ; 30 cm.\\
	
	\hspace{0.5cm} \imprimirorientadorRotulo~\imprimirorientador\\
	
	\hspace{0.5cm}
	\parbox[t]{\textwidth}{\imprimirtipotrabalho~--~\imprimirinstituicao,
	\imprimirdata.}\\
	
	\hspace{0.5cm}
		1. Palavra-chave1.
		2. Palavra-chave2.
		I. Orientador.
		II. Universidade xxx.
		III. Faculdade de xxx.
		IV. Título\\ 			
	
	\hspace{8.75cm} CDU 02:141:005.7\\
	
	\end{minipage}
	\end{center}
	\hrule
\end{fichacatalografica}
\end{comment}
% ---


% ---
% Inserir Errata (opcional)
% ---
\begin{comment}
\begin{errata}
Elemento opcional da \citeonline[4.2.1.2]{NBR14724:2011}. Exemplo:

\vspace{\onelineskip}

FERRIGNO, C. R. A. \textbf{Tratamento de neoplasias ósseas apendiculares com
reimplantação de enxerto ósseo autólogo autoclavado associado ao plasma
rico em plaquetas}: estudo crítico na cirurgia de preservação de membro em
cães. 2011. 128 f. Tese (Livre-Docência) - Faculdade de Medicina Veterinária e
Zootecnia, Universidade de São Paulo, São Paulo, 2011.

\begin{table}[htb]
\center
\footnotesize
\begin{tabular}{|p{1.4cm}|p{1cm}|p{3cm}|p{3cm}|}
  \hline
   \textbf{Folha} & \textbf{Linha}  & \textbf{Onde se lê}  & \textbf{Leia-se}  \\
    \hline
    1 & 10 & auto-conclavo & autoconclavo\\
   \hline
\end{tabular}
\end{table}

\end{errata}
\end{comment}
% ---

% ---
% Inserir Folha de aprovação (obrigatório)
% ---

% Isto é um exemplo de Folha de aprovação, elemento obrigatório da NBR
% 14724/2011 (seção 4.2.1.3). Você pode utilizar este modelo até a aprovação
% do trabalho. Após isso, substitua todo o conteúdo deste arquivo por uma
% imagem da página assinada pela banca com o comando abaixo:
%
%\includepdf[pages=-]{folhadeaprovacao.pdf}
%

%TODO Adicionar folha de aprovação
\begin{folhadeaprovacao}

\includepdf[scale=0.9, pages=1-2, pagecommand={}]{folhadeaprovacao.pdf}
  
\end{folhadeaprovacao}
% ---

% ---
% Dedicatória (opcional)
% ---
% \begin{dedicatoria}
%    \vspace*{\fill}
%    \centering
%    \noindent
%    \textit{Dedico este trabalho à minha família...} \vspace*{\fill}
% \end{dedicatoria}
% ---

% ---
% Agradecimentos (opcional)
% ---
\begin{agradecimentos}

%TODO
[TO DO] Escreva seu agradecimento

\end{agradecimentos}
% ---

% ---
% Epígrafe (opcional)
% ---
\begin{epigrafe}
    \vspace*{\fill}
	\begin{flushright}
		\textit{"A persistência é o caminho do êxito." (Charles Chaplin)}
	\end{flushright}
\end{epigrafe}
% ---

% ---
% RESUMOS
% ---

% resumo em português
% Resumo em língua vernácula (obrigatório)
\begin{resumo}
% TODO: escrever resumo do TCC (150-300 palavras) cobrindo: contexto, problema, objetivo, metodologia, resultados (quando houver), conclusao.

\vspace{\onelineskip}
\noindent
\textbf{Palavras-chaves}: poda neural; eficiência energética; compressão de modelo de IA; otimização estrutural de redes neurais.
\end{resumo}

% resumo em inglês
% Resumo em língua estrangeira (obrigatório)
\begin{resumo}[\textbf{Abstract}]
 \begin{otherlanguage*}{english}

[TO DO] Abstract em um parágrafo.
    
   \vspace{\onelineskip}
 
   \noindent 
   \textbf{Key-words}: [TO DO] Key-words separadas por ponto-e-vírgula
 \end{otherlanguage*}
\end{resumo}
% ---

% ---
% inserir lista de ilustrações  essa lista é prenchida automatica  com as imagens inseridas no documento
% Listas de ilustrações (opcional)
% ---
\pdfbookmark[0]{\listfigurename}{lof}
\listoffigures*
\cleardoublepage
% ---

% ---
% inserir lista de tabelas é prenchida automatica  com as tabelas  inseridas no documento
% Lista de tabelas (opcional)
% ---
\pdfbookmark[0]{\listtablename}{lot}
\listoftables*
\cleardoublepage
% ---

% ---
% inserir lista de abreviaturas e siglas inserir caso houver
% Lista de abreviaturas (opcional)
% Lista de siglas (opcional)
% ---
\begin{siglas}
    \item [TO DO] {\em [TO DO]}
 \end{siglas}
% ---


% ---
% inserir lista de símbolos
% Lista de símbolos (opcional)
% ---
\begin{comment}
\begin{simbolos}
  \item[$ \Gamma $] Letra grega Gama
  \item[$ \Lambda $] Lambda
  \item[$ \zeta $] Letra grega minúscula zeta
  \item[$ \in $] Pertence
\end{simbolos}
\end{comment}
% ---

% ---
% inserir lista de algoritmos lista os algoritimos automatico
% Lista de algoritmos (opcional)
% ---
\pdfbookmark[0]{\listtablename}{loa}
\listofcodes
\cleardoublepage
% ---


% ---
% inserir o sumario é preenchido automatico de acordo que é criado no documento
% Sumário (obrigatório)
% ---
\pdfbookmark[0]{\contentsname}{toc}
\textbf{\tableofcontents*}
\cleardoublepage
% ---